\section{Conclusion and Future Directions}

This paper studies online scheduling for LLM inference, a setting in which each request's memory consumption grows during decoding and exceeding capacity triggers evictions that can destabilize a system whose capacity is theoretically sufficient for stability. We develop a multi-stage online scheduling model that captures these dynamics and analyze a fluid approximation whose equilibrium identifies the stability region and the memory level $M^*$ required for stability at its boundary. Guided by this fluid benchmark, we design the (Nested) WAIT algorithm, a threshold-based admission rule that keeps the stochastic system close to the fluid equilibrium. The threshold mechanism performs a dimensionality reduction: it decouples interactions across prompt types and replaces the continuously growing KV cache dynamics with a discrete-time per-type recursion. Under this reduction, we establish asymptotic optimality for the effective throughput, the completion rate net of eviction, even when output lengths are unknown at arrival. Because the algorithms operate near the fluid equilibrium and prevent the eviction cascades that destabilize memory-reactive baselines, they both enlarge the stability region and reduce mean end-to-end latency across the underloaded, near-capacity, and overloaded regimes. Numerical experiments on Llama-2-7B with an NVIDIA A100 GPU confirm both effects relative to vLLM and Sarathi.

Several questions remain open. First, our threshold design relies on knowledge of the prompt-type arrival distribution; developing distribution-free variants, or versions that adapt online under bounded misspecification, would extend the algorithm's applicability. Second, practical deployments face arrival rates that drift over time and occasionally surge. Theorem~\ref{thm:nested_wait_heavy_traffic_time_varying} handles bounded rate variation by periodically re-solving the fluid equilibrium, but sharper volatility, such as demand spikes that temporarily exceed nominal capacity, calls for further analysis; two concrete questions are how fast the thresholds can re-tune while preserving asymptotic optimality, and how transient surges interact with the safety buffer in Nested WAIT.

Finally, extending the analysis to multi-GPU deployments introduces inter-device communication cost and a choice of parallelization strategy (data, tensor, or pipeline). This regime is especially relevant for mixture-of-experts (MoE) architectures \citep{shazeer2017outrageously, fedus2022switch, jiang2024mixtral}, in which token-level routing distributes experts across GPUs and creates skewed memory and compute load that our single-GPU analysis does not capture. A unified framework that models inter-device communication inside the iteration-time model, and coordinates admission thresholds across GPUs in response to the routing distribution, is a natural open direction.
